% Generated draft: requires mathematical and source review.
\documentclass[11pt]{article}
\usepackage[margin=25mm]{geometry}
\usepackage{fontspec}
\setmainfont{Latin Modern Roman}
\usepackage{amsmath,amssymb,mathtools,unicode-math}
\setmathfont{Latin Modern Math}
\usepackage{enumitem,needspace,xurl,hyperref}
\setlength{\parindent}{0pt}
\setlength{\parskip}{5pt}
\setlength{\emergencystretch}{3em}
\newcommand{\N}{\mathbb{N}}
\newcommand{\Z}{\mathbb{Z}}
\newcommand{\Q}{\mathbb{Q}}
\newcommand{\R}{\mathbb{R}}
\newcommand{\C}{\mathbb{C}}
\newcommand{\F}{\mathbb{F}}
\newcommand{\A}{\mathbb{A}}
\newcommand{\PP}{\mathbb{P}}
\newcommand{\E}{\mathbb{E}}
\newcommand{\eps}{\varepsilon}
\newcommand{\dd}{\,\mathrm{d}}
\begin{document}
\section*{UnsolvedMath ID: 20000426}
\textbf{Source problem code:} catalogue entry
\subsection{Definitions and mathematical statement}
Let \(X\) be a modular curve parametrizing elliptic curves with a specified level structure, and let \(\mathcal E\to X\) be its universal family where such a family exists. Choose an appropriate smooth projective elliptic-surface model \(S\), typically over \(\Q\), and a prime of good reduction with special fiber \(S_p\). For \(\ell\ne p\), a divisor Tate class is a class in \(H^2_{\mathrm{\acute et}}((S_p)_{\overline{\F}_p},\Q_\ell(1))\) fixed after a finite extension of the residue field. The problem attributed to J. Ellenberg and D. Ulmer asks for algebraic cycles on characteristic-\(p\) fibers representing the new Tate classes not obtained by specialization from characteristic zero. It is a construction request in the primes and surfaces where such new classes occur, not an assertion that every level and every prime has a jump. The level, integral model and precise comparison/specialization convention are not fixed in the original sentence.

The supplied partial result treats only the full level-4 elliptic modular surface at \(p=3\), under an identification of its reduction with a Fermat quartic. Write \(S_0\) for the geometric characteristic-zero surface, \(S_3\) for the geometric special fiber, and \(\operatorname{sp}:\operatorname{NS}(S_0)\to\operatorname{NS}(S_3)\) for its N\'eron--Severi specialization. The N\'eron--Severi group is the group of divisor classes modulo algebraic equivalence. In the report's Fermat coordinates over \(\F_9=\F_3(i)\), \(i^2=-1\), the two proposed lines are
\[
 L_2:\quad x_1+(-1+i)x_4=x_2+(-1-i)x_3=0,
\]
\[
 L_{19}:\quad x_1+x_3+x_4=x_2+i x_3-i x_4=0.
\]
The candidate contribution asserts
\[
 \operatorname{NS}(S_3)=\operatorname{sp}(\operatorname{NS}(S_0))
 \oplus\Z[L_2]\oplus\Z[L_{19}].
\]
This is an integral direct sum of abelian groups, not an assertion of orthogonality for the intersection pairing. The report states that adjoining these lines to the twenty specialization rows in Shimada's published integral bases gives determinant \(-1\), providing a minimal unimodular integral complement and accounting for the full two-dimensional geometric Tate-class jump.

It reports a quotient Frobenius matrix
\[
 \begin{pmatrix}1&1\\0&-1\end{pmatrix},
\]
with one new invariant direction over \(\F_3\), and both new directions Tate over \(\F_9\). The integral orbit sum \(L_2+L_5\) is reported to be an \(\F_3\)-defined algebraic cycle spanning that new rational invariant direction. The label \(L_5\), the Frobenius basis/action convention, the model identification, the twenty rows and the integral bases are not defined explicitly in the supplied record. Thus the displayed construction is retained as a reported special-case result requiring those verifications, not a general construction for all modular elliptic surfaces or a claim to have reproduced its determinant calculation.
\subsection{Short English statement}
Find characteristic-\(p\) cycles accounting for Tate classes newly appearing in reductions of modular elliptic surfaces. The supplied report proposes two explicit \(\F_9\)-lines for the level-4 surface modulo three and an \(\F_3\)-defined orbit sum, while leaving the general problem open and the basis/model computations for review.
\subsection{Sources}
\begin{itemize}
\item Local UnsolvedMath registry, permanent ID \(20000426\), source code \texttt{catalogue entry}; AIM algebraic number theory corpus, \texttt{aim-algebraic-number-theory-notes.json}, notes[77], stored source item 5. Attribution: Ulam AI / UnsolvedMath Contributors, CC BY 4.0.
\item \emph{Problems from the AIM Tate Conjecture Workshop}, July 23--27, 2007, transcribed by Christopher Lyons, attributed to J. Ellenberg and D. Ulmer: \url{https://aimath.org/WWN/tateconjecture/tateconjecture.pdf}. The supplied verification reports actual Problem 15, correcting OCR ``5'', the misspelling ``paramterizes'', and a stray page marker ``1''; both numbering forms are retained as source information.
\item The research summary cites Shimada's published integral bases without title, date, full reference, or the basis and specialization matrices. No missing bibliographic details or equation for \(L_5\) are invented.
\item Registry dated August 14, 2026: statement exact, status \texttt{partially\_solved}, classification \texttt{PARTIAL-PROGRESS}; attempt 1 records problem status open, partial result, rigor proved, novelty confidence low. The explicit-cycle and Frobenius claims are machine-generated research aids, not peer reviewed, requiring independent verification before citation.
\end{itemize}
\end{document}
